% Kurzbeschreibung intelligente_Suchabfrage_7.py
% Maximal 30 Zeilen, Stichpunkte / kurze Sätze

\textbf{intelligente\_Suchabfrage\_7.py -- Freitextsuche nach Kandidaten}
Nimmt eine natürlichsprachliche Suchanfrage (z.B. \glqq Facharzt Radiologie in München\grqq{}) entgegen.
Extrahiert aus der Anfrage Intent-Informationen: Position, Fachbereich, Stadt, Bundesland und Land.
Nutzt vordefinierte Listen für Positionen, Fachrichtungen und Länder (inkl. Abkürzungen wie \glqq DE\grqq{}).
Lädt alle Kandidaten aus der Tabelle \texttt{candidates} und normalisiert zentrale Textfelder.
Interpretiert \texttt{regionale\_verfuegbarkeit} analog zum Matching (deutschlandweit vs. km-Radius).
Liest Städte- und Bundeslandinformationen aus der Datei \texttt{Staedte\_Deutschland.csv} ein.
Berechnet mit der Haversine-Formel Entfernungen zwischen Suchort und Wohn- bzw. Wunscharbeitsort.
Prüft optional Karrierestufen-Logik (Position des Kandidaten vs. gesuchte Position).
Prüft optional den Fachbereich des Kandidaten gegen den gewünschten Fachbereich.
Kandidaten werden zuerst über maximale Entfernung (z.B. 50 km) gefiltert.
Falls Entfernung zu groß ist, wird geprüft, ob Wohn- oder Wunscharbeitsort im gleichen Bundesland liegt.
Falls kein Bundesland-Treffer, wird geprüft, ob der Wunscharbeitsort im gleichen Land liegt (inkl. Kürzel wie \glqq de\grqq{}).
Kandidaten mit Wunscharbeitsort \glqq deutschlandweit\grqq{} bzw. nur \glqq DE\grqq{} werden immer berücksichtigt.
Erstellt aus allen passenden Kandidaten eine kompakte Ergebnisliste mit Entfernung, Wohnort und Wunscharbeitsort.
Sortiert die Ergebnisliste nach aufsteigender Entfernung (kürzester Fahrtweg zuerst).
Kann alle Rohdaten der passenden Kandidaten als Excel-Datei im Ordner \texttt{results} exportieren.

